\chapter{Literature Review}

Some of the research carried out on Cross-Dataset Generalization for NIDS have shown improvements with advanced dataset handling and feature engineering 
on supervised machine learning and deep learning models. These models have achieved greater accuracy on benchmark datasets and improved to some extent 
in cross-dataset generalization (between CICIDS2017 and CSE-CIC-IDS2018, UNSW-NB15). However, accuracy alone cannot be considered to determine how robust 
the model is given the high number of genuine network flows and a smaller number of attacks. Cross-Dataset Generalization for NIDS face reliability issue 
when models are deployed for intrusion detection considering the increasing number of attacks in heterogeneous environments with huge domain shifts when 
the models with greater accuracy on benchmark dataset fail to detect intrusion in an unknown dataset yielding more False Positives.  Generalization becomes 
complex as the features from the train dataset that are predictive of malicious activities within the dataset may fail to predict malicious behavior from a 
different dataset thus dropping the performance of the models \cite{Elangovan2026}. Uniformity in dataset size for training and testing, usage of SMOTE 
are crucial to make comparisons on how the models performed in intrusion detection in the future. 

The studies have shown the effectiveness of ensemble learning and feature selection in finding nonlinear dependencies among different labels in the 
network flow by keeping low dimensionality. They have also cautioned to interpret the evaluation wisely despite Random Forest, XGBoost, 1D-CNN achieving 
accuracy \>85\% on benchmark datasets \cite{Elangovan2026}, these models faced challenges modelling on complex multi-stage attacks \cite{Chandra2025}. 

\section{Cross-Dataset Generalization using Supervised Learning Models}
With the ever-increasing study and experiments on network intrusion detection, a recent study by \cite{Cantone2024}, although not validated and reviewed
 by experts performed cross-dataset generalization by employing machine learning algorithms on CIC-IDS-2017, CSE-CIC-IDS2018, LycoS-IDS2017, 
 and LycoS-UnicasIDS2018 datasets and eventually applied corrections to CSE-CIC-IDS2018 based on the results of LycoS-IDS2017. The study revealed that the
  performance dropped to detect malicious behaviors for generalization in comparison to near perfect results in intrusion detection when trained and tested 
  on same datasets. The supervised learning models learned bigger scale attack types such as DOS and DDoS and generalized well. Meanwhile, due to class imbalance
   of the attack types, the models failed to detect web attacks and infiltration attacks. The researchers have pointed out that the variability of the datasets, 
   irregular distribution of attacks among datasets and redundancy in the samples played a causal role in performance drop of the models failing to learn the 
   attack patterns. Problems associated with feature calculation, protocol identification, feature multiplication and labeling in CIC-IDS-2017 degraded the 
   quality of the data and hence corrected with LycoS-IDS2017. The findings suggested that the learning of the transferrable attacks might not be the 
   real reason, but the features associated with the dataset that hinder the performance of the models. As per the study, the models should be implemented 
   with prior knowledge of different traffics that were never encountered before. 

The table below showcases the performance of the models while trained on CIC-IDS-2017 and tested on CSE-CIC-IDS2018 as per \cite{Cantone2024} for 
cross-dataset generalization for network intrusion detection.

\begin{table}[ht]
    \centering
    \caption{ Model Evaluation Results from \cite{Cantone2024}.}
    \label{tab:evaluation-results}

    \begin{tabular}{ll l c c c}
        \toprule
        {\small Train Dataset} & {\small Test Dataset} & {\small Models} &
        \multicolumn{3}{c}{\small Evaluation Metrics} \\
        \cmidrule(lr){4-6}
        & & & \textbf{{\small MCC}} & \textbf{{\small F1}} & \textbf{{\small ROC-AUC}} \\
        \midrule

        \multirow{3}{*}{\textbf{\small CIC-IDS-2017}}
        & \multirow{3}{*}{\textbf{\small CSE-CIC-IDS2018}}
        & {\small Random Forest} & {\small 25.96\%} & {\small 25.76\%} & {\small 66.93\%} \\

        & & {\small XGBoost} & {\small 35.68\%} & {\small 26.43\%} & {\small 80.44\%} \\

        \bottomrule
    \end{tabular}
\end{table}
\FloatBarrier

Another study by \cite{Elangovan2026} revealed the performance drop for NIDS while evaluated on recent datasets. Random Forest when trained 
on CIC-IDS-2017 and tested on CSE-CIC-IDS2018 saw a Macro F1 drop of 0.059 from baseline 0.941 to 0.882. Similarly, 1D-CNN saw a decline of 0.070 
from 0.942 to 0.872. It seems obvious that the performance degrades while generalizing to a completely different dataset when both temporal and 
network domain distance increases implying that performance on benchmark dataset does not signify robustness for future attack detection. The 
study highlights the importance of carrying out temporal forecasting evaluation while implementing models for intrusion detection.  


\section{Hybrid Deep Learning approach for Cross-Dataset Generalization}
In addition to the application of supervised learning models, usage of deep learning methods has been observed for the detection of intrusion due to 
its capability of learning hierarchical representations from the complex traffic data. A research paper “Real-Time Network Intrusion Detection System 
Using Deep Learning: Identifying Cyber Attacks Through Anomaly Pattern Recognition” \cite{Chandra2025} has properly analyzed the performance of LSTM, 
CNN and Autoencoders in detecting anomalies. The paper has discussed the hybrid approach and proposed the Real-Time NIDS by employing attention mechanism, LSTM, 
CNN and Autoencoders. In this hybrid architecture setup, LSTMs model network flow sequences, CNNs identify spatial patterns from the traffic features, Autoencoder 
calculates the anomaly score from the reconstruction error and lastly the important characteristics/temporal phases within the traffic data sequence are identified 
by attention mechanism.

Like the performance of the supervised learning models on benchmark datasets, application of deep learning for intrusion detection recorded exemplary F1 scores 
of 97.8\%, 97.3\% and 98.0\% for CIC-IDS-2017, CSE-CIC-IDS2018 and UNSW-NB15 dataset respectively. The study has conducted cross-dataset performance between CIC-IDS-2017 
and UNSW-NB15 only. Using this hybrid approach, they have yielded a F1-score of 93.5\% and FPR of 2.8\% only which has surpassed the performance of the base supervised 
learning models. The major difference in the results between the application of supervised learning models from previous research and the hybrid DNN approach lie in 
data split for training and testing. The study of cross-dataset generalization by \cite{Cantone2024} used whole dataset for training and testing. Meanwhile, 
the research paper by \cite{Chandra2025} used SMOTE (to oversample rare attack classes), 70\% of each dataset for training, 15\% for hyperparameter tuning and remaining 15\% 
for testing. While this approach yielded superior performance, as per Elangovan et al. (2026), this complex deep learning architecture does not guarantee a reliable 
generalization as the performance for intrusion detection degrade when there are large distribution shifts and may not be resilient enough to identify unfamiliar traffic 
thus contradicting the notion that complex model would yield better performance when the models are trained on older traffics and tested on newer version of web traffics 
with evolving network technologies, communication behaviors and varied attack patterns. Despite the robustness of the models, \cite{Elangovan2026} has stated the necessity 
of preserving the network features when the network traffic features are classified based on their semantics during dataset shift. Protocol specific indicators showing 
higher degree of uncertainty whereas communication features like flow rate and temporal activity features generalizing well on different datasets. The negative communication 
behavior aspects that arise in different situations can be handled by certain behavioral characteristics. The evidence for this is discussed in the representation-drift 
study by \cite{Elangovan2026}. Higher the difference between training and testing datasets with respect to time and field, higher the construction error. Models trained 
on CIC-IDS-2017 showed 0.0211 mean reconstruction error while analyzed on CSE-CIC-IDS2018, increased to 0.0342 for CICIoT2023 and reached the maximum of 0.0384 when 
applied to CIC-DIAD2024 showing the latent representations obtained from the original dataset diverged from the target-domain traffic.

The findings from the research for cross-dataset generalization by \cite{Chandra2025,Elangovan2026,Cantone2024} are in contrast considering the models used. 
These papers show how the evaluation design, datasets, pre-processing, feature engineering and model design affect generalization. Thus, cross-dataset generalization for 
NIDS remains continuous deep research of the specific source and target datasets, features compatibility, attack overlap and experimental approach.

\section{Research Gap}
All these previous studies of Cross-Dataset generalization for NIDS using Supervised Machine Learning and Deep Learning Models demonstrate gradual 
progress for NIDS, although the evaluations were made using different metrics on different test datasets. Trees based models seem to provide faster classification 
benchmarks while deep learning methods contribute to representation learning, temporal modelling, and anomaly detection. Hybrid deep learning method 
could significantly assist real-time cross-dataset generalization for NIDS. However, achieving accuracy greater than 80\% or even 90\% with high 
precision scores on a single benchmark dataset or a cross-dataset (as the results suggested by the studies above) does not confirm its robustness 
for deployment. They fail to exhibit excellent generalization capabilities. There is no consistent performance across varied test datasets thus failing
 to deliver similar results. The inconsistency in performance across datasets proves the significance of good labeling, dataset preparation and feature 
 engineering. There remains a consistency gap in the advancement of NIDS evaluation approaches that would consider the efficiency of classification and 
 generalization to different datasets, robustness over time, feature stability, imbalance class, false-positive estimations, and computational efficiency. 
 Future research should conduct collaborative experiments for real deployment simulations with independent datasets and heterogeneous network conditions.

 \section{Research Questions}
To serve as a guide for this research overcoming the research gaps discussed earlier, we have developed the 'Qualitative' research questions that specifically
contribute towards the efficiency of both supervised and unsupervised machine learning models for network intrusion detection on different
datasets.

The following are the qualitative research questions developed for the research:

\begin{enumerate}
    \item How well do supervised machine learning models trained on CIC-IDS-2017 dataset generalize to unseen intrusion detection datasets?
    \item How does the cross-dataset generalization performance of unsupervised models compare with supervised models?
    \item Which machine learning model demonstrate greater robustness against dataset distribution shift in real-world network intrusion detection scenarios?
\end{enumerate}

 \section{Research Hypotheses}
 Cross-dataset generalization for network intrusion detcetion is crucial in determining the robustness of the models in identifying intusion from a varied 
 unknown network traffic. As previous studies have stated that although the models demonstrated exemplary performance while trained and tested on same models;
 the models saw significant drop in performance and struggled to generalize efficiently with respect to unseen dataset environment comprising of varied traffic
 characteristics and attack types. Based on these facts, we have formulated few hypotheses to investigate the cross-dataset generalization capabilities of the 
 supervised and unsupervised learning models for developing a robust network intrusion detection system.

 \begin{itemize}
    \item \textbf{H1:} While trained on a single dataset and tested on other unseen datasets, cross-dataset generalization may not be significant despite higher accuracy achieved by the models.
    \item \textbf{H2:} Despite the models generating higher Precision score, there will be a reduced or unsignificant Recall score during generalization. 
    \item \textbf{H3:} Supervised machine learning models will yield significantly better results than Unsupervised machine learning models.
    \item \textbf{H4:} Tuning hyperparameters alone will not be enough for DNN to be robust in identifying intrusions. Number of trials should be increased too to learn the patterns. 
    \item \textbf{H5:} The presence of completely different attack types in UNSW-NB15 dataset, models trained on CIC-IDS-2017 and tested on UNSW-NB15 will see performance decline
    compared to models tested on CSE-CIC-IDS2018.
 \end{itemize}


 \section{Research Objectives}
The ultimate goal of this research is to examine the performances of supervised and unsupervised machine learning models, evaluate their capabilities in 
cross-dataset generalization for network intrusion detection. The research follows an experiment by training the models on CIC-IDS-2017 and evaluate their 
performance in CSE-CIC-IDS2018 and UNSW-NB15 datasets to identify different types of network attacks. We have set the following objectives to achieve our goal.

\begin{enumerate}
    \item Process and standardize the CIC-IDS-2017, CSE-CIC-IDS2018 and UNSW-NB15 datasets.
    \item Train supervised and unsupervised machine learning models.
    \item Evaluate the model performance on CSE-CIC-IDS2018 and UNSW-NB15 test datasets.
    \item Make comparison of the evaluation metrics: F1\_Score, MCC, ROC-AUC, Precision, and Recall.
    \item Identify and select the best model that is more robust to dataset shift and generalizes better for intrusion detection in real-world cybersecurity environments.
\end{enumerate}
